% Authority: exact-scope leaf 26; full PDF page 105; printed p.88. Direct transcription.
% U:p088.u01 | continuation
\noindent
Cette dernière congruence fait voir que $-b$ est ou n'est pas résidu biquadratique
relativement à $p$, selon que l'on a $\leg{t}{p}=1$, ou $\leg{t}{p}=-1$. Si l'on compare
ce résultat à la formule $(\gamma')$, on obtiendra cet énoncé:
% U:p088.u02 | labelled-statement
\begin{equation*}\tag{$\delta'$}
\begin{minipage}{0.83\textwidth}
$-b$ est ou n'est pas résidu biquadratique par rapport à $p$, selon
que l'on a $\leg{t}{b}=1$ ou $\leg{t}{b}=-1$.
\end{minipage}
\end{equation*}
% U:p088.u03 | paragraph-start
\par\indent
Décomposons actuellement le nombre premier $p$ en deux carrés, en posant
$p=\varphi^2+\psi^2$ (où $\psi$ est supposé pair) et mettons cette valeur de $p$ dans
l'équation $(\alpha')$. Il viendra ainsi:
\[
t^2-bu^2=s^2\varphi^2+s^2\psi^2,
\]
et l'on aura ensuite en transposant:
\[
t^2-\psi^2s^2=(t+\psi s)(t-\psi s)=s^2\varphi^2+bu^2.
\]
% U:p088.u04 | paragraph-start
\par\indent
Les nombres $s$, $t$, $u$ étant deux à deux premiers entre eux, on voit
par l'équation $(\alpha')$ que $s$ ne saurait être divisible par $b$. Nous distinguerons
maintenant deux cas selon que $\varphi$ est ou n'est pas divisible par $b$, et nous
examinerons d'abord le dernier de ces deux cas. Si nous désignons par
$m$ le plus grand diviseur commun de $\varphi$ et $u$, $m$ sera impair et non-divisible
par $b$ comme $\varphi$. Posons $\varphi=m\varphi'$, $u=mu'$ et substituons ces valeurs dans
la dernière équation. Il viendra ainsi:
\[
(t+\psi s)(t-\psi s)=m^2(s^2\varphi'^2+bu'^2).
\]
Comme $s\varphi'$ et $bu'$ sont premiers entre eux, d'après ce qui précède, il suit
d'un théorème connu, conséquence très simple de la loi de réciprocité (\emph{Théorie
des Nombres} no. 197), que tout diviseur $R$ du facteur binôme du second membre
qui est évidemment impair, est tel que $\leg{R}{b}=1$.
% U:p088.u05 | paragraph-start
\par\indent
La dernière équation fait voir que chacun des nombres $t+\psi s$, $t-\psi s$
est composé de deux facteurs, l'un diviseur de $m^2$, l'autre diviseur de $s^2\varphi'^2+bu'^2$.
Si nous désignons ces 4 facteurs par $E$, $F$, $K$ et $L$, nous aurons les équations:
\[
\begin{array}{ll}
t+\psi s=EK,&m^2=EF,\\
t-\psi s=FL,&\varphi'^2s^2+bu'^2=KL.
\end{array}
\]
% U:p088.u06 | paragraph-start; continues on p089
\par\indent
Il est facile de s'assurer que les nombres $E$ et $F$, qui sont impairs l'un
et l'autre, comme diviseurs de $m^2$, sont premiers entre eux. En effet, si ces
nombres étaient divisibles l'un et l'autre par un même nombre premier $\delta$, $\delta$
diviserait chacun des nombres $t+\psi s$, $t-\psi s$ et par conséquent aussi leur demi-somme
et leur demi-différence, c'est-à-dire les nombres $t$ et $\psi s$. D'un autre
